\Needspace{13\baselineskip}
\section{Divisorial constraints and residual volume}
\label{sec:divisorial-gaud-volume}

Suppose we require a $\theta$-plurisubharmonic function to have generic Lelong numbers at least $a_j$ along prescribed prime divisors $D_j$. Removing the divisor $D=\sum_ja_jD_j$ leaves the residual background $\theta-\eta_D$, where $\eta_D$ is a curvature form of $D$, and in each Gauduchon quotient this removal subtracts the Gauduchon degree of $D$ from the numerator. Lamari's duality therefore decides when such functions exist, and the volume of the residual background measures what remains. By the correspondence of \cref{thm:ordinary-divisor-normalization}, equations in the ordinary singularity type of the logarithmic potential of $D$ reduce to equations in the residual background. When the residual background has a smooth positive representative, the smooth maximizer of the dual problem forces every subsolution with the prescribed Lelong numbers and the residual volume as scalar to be the logarithmic potential of $D$ plus a smooth function, up to a constant.

Fix a smooth real $(1,1)$-form $\theta$ on $X$. Let $D=\sum_{j=1}^N a_jD_j$ be an effective real divisor, with distinct prime divisors $D_j$ and coefficients $a_j>0$. Choose a smooth Hermitian metric on each line bundle $\mathcal O_X(D_j)$, write $s_j$ for its canonical section, and let $\eta_j$ be its Chern curvature form. Set
\[
 \sigma_D=\sum_j a_j\log|s_j|^2,
 \qquad \eta_D=\sum_j a_j\eta_j,
 \qquad \beta_D=\theta-\eta_D.
\]
The Poincar\'e--Lelong formula reads
\[
 \ddc\sigma_D=[D]-\eta_D.
\]
For a Gauduchon metric $\gamma$ on $X$, define the Gauduchon degree
\[
 \deg_\gamma D=\int_X[D]\wedge\gamma^{n-1}
              =\sum_j a_j\int_{D_j}\gamma^{n-1},
\]
where integration over a singular divisor means integration over its regular locus, or equivalently against its current of integration. This degree is finite, positive for $D\ne0$, and equal to $\int_X\eta_D\wedge\gamma^{n-1}$ by Stokes' theorem for currents. We write $\nu(u,D_j)$ for the generic Lelong number of $u$ along $D_j$, as in \cref{sec:analytic-types}.

\Needspace{12\baselineskip}
\begin{proposition}\label{prop:divisorial-gaud-feasibility}
The following are equivalent:
\begin{enumerate}
\item there is $u\in\PSH(X,\theta)$ with $\nu(u,D_j)\geq a_j$ for every $j$;
\item the residual background $\beta_D$ is pseudoeffective;
\item for every Gauduchon metric $\gamma$,
\[
 \int_X\theta\wedge\gamma^{n-1}\geq\deg_\gamma D.
\]
\end{enumerate}
The first condition does not require $u-\sigma_D$ to be bounded, nor the Lelong numbers of $u$ to equal the prescribed ones.
\end{proposition}
\begin{proof}
Suppose first that $u$ satisfies the first condition. By \cref{lem:nonclosed-divisorial-removal}, these Lelong numbers can be subtracted: $\theta+\ddc u-[D]\geq0$. Consequently
\[
 \beta_D+\ddc(u-\sigma_D)=\theta+\ddc u-[D]\geq0,
\]
where the difference $u-\sigma_D$ is first taken in $L^1$ and then replaced by its quasi-plurisubharmonic representative. This proves the second condition. Conversely, if $v\in\PSH(X,\beta_D)$, then $u=\sigma_D+v$ satisfies
\[
 \theta+\ddc u=[D]+\beta_D+\ddc v\geq0,
\]
so its generic Lelong numbers along the $D_j$ are at least $a_j$. For $n\geq2$, Lamari's duality, \cref{thm:lamari}, says that a smooth background is pseudoeffective exactly when all its Gauduchon integrals are nonnegative; applied to $\beta_D$, together with the degree identity, it gives the equivalence of the second and third conditions. For $n=1$ the same criterion reads $\int_X\beta_D\geq0$: necessity follows from Stokes' theorem, and for sufficiency one solves $\ddc f=(\int_X\beta_D)\mu-\beta_D$ for a smooth real function $f$ with the scalar Green operator.
\end{proof}

For a nonzero effective real divisor $D$ and a pseudoeffective $\theta$, the exact threshold is therefore
\begin{equation}\label{eq:divisorial-gaud-threshold}
 \tau_\theta(D)=
 \inf_{\gamma\text{ Gauduchon}}
       \frac{\int_X\theta\wedge\gamma^{n-1}}{\deg_\gamma D}.
\end{equation}
It is finite and nonnegative, and a function $u\in\PSH(X,\theta)$ with $\nu(u,D_j)\geq ta_j$ for all $j$ exists if and only if $0\leq t\leq\tau_\theta(D)$. The endpoint is included, because the non-strict Gauduchon inequalities still hold there and duality gives a positive current. This is a threshold for the existence of a positive current, not for the existence of a bounded residual potential.

Whenever $\beta_D$ is pseudoeffective, define the residual volume
\[
 \mathcal V_\mu(\theta;D)=\mathcal V_\mu(\beta_D).
\]
The definition and the degree identity give the exact formula
\begin{equation}\label{eq:divisorial-gaud-volume}
 \mathcal V_\mu(\theta;D)^{1/n}
   =\inf_{\gamma\text{ Gauduchon}}
     \frac{\int_X\theta\wedge\gamma^{n-1}-\deg_\gamma D}
          {D_\mu(\gamma)},
\end{equation}
and the residual volume is finite and nonnegative. Changing the auxiliary metrics on the line bundles changes $\eta_D$ by a smooth $\ddc$-exact term, so the residual volume does not depend on these metrics, nor on a smooth $\ddc$ change of $\theta$. It depends on the prescribed divisor and on the smooth probability measure $\mu$.

Here $\mathrm{MA}^{\mathrm{ext}}_\theta$ denotes the extended non-pluripolar Monge--Amp\`ere measure of \cref{thm:arbitrary-extended-ma}.

\begin{theorem}\label{thm:divisorial-residual-equation}\index{Divisor!residual equation}
Suppose that $\beta_D$ has a smooth positive representative, and put $c=\mathcal V_\mu(\theta;D)$. There is a $\theta$-plurisubharmonic function $u_0=\sigma_D+v$, with $v$ smooth, satisfying
\[
 \mathrm{MA}^{\mathrm{ext}}_\theta(u_0)=c\mu.
\]
Every $u\in\PSH(X,\theta)$ such that
\[
 \nu(u,D_j)\geq a_j\quad(1\leq j\leq N),\qquad
 \mathrm{MA}^{\mathrm{ext}}_\theta(u)\geq c\mu
\]
differs from $u_0$ by a constant. In particular, a solution of a normalized equation $\mathrm{MA}^{\mathrm{ext}}_\theta(u)=c'\mu$ in the ordinary singularity type of $\sigma_D$ has scalar $c'=c$, has a smooth remainder $u-\sigma_D$, and is unique up to constants.
\end{theorem}
\begin{proof}
\emph{Existence.} The smooth equation and \cref{thm:gaud-volume-dual-formula} give a smooth $v$ with
\[
 \Omega=\beta_D+\ddc v>0,\qquad \Omega^n=c\mu.
\]
Then $\theta+\ddc(\sigma_D+v)=[D]+\Omega\geq0$. Off $D$, the extended measure of $u_0=\sigma_D+v$ is $\Omega^n$, and both measures give zero mass to $D$. This proves existence.

\emph{Reduction to the residual background.} Let $u$ satisfy the two inequalities. Removing the divisor as in the proof of \cref{prop:divisorial-gaud-feasibility}, we obtain
\[
 w=u-\sigma_D-v\in\PSH(X,\Omega),\qquad
 \Omega+\ddc w=\theta+\ddc u-[D]\geq0,
\]
where the difference is first taken in $L^1$ and then replaced by its canonical representative. On $X\setminus D$, the two currents differ only by smooth changes in the decomposition of their local potentials, so their extended measures agree there, and both measures vanish on $D$. Hence
\[
 \mathrm{MA}^{\mathrm{ext}}_\Omega(w)
       =\mathrm{MA}^{\mathrm{ext}}_\theta(u)\geq c\mu.
\]

\emph{Rigidity.} Take $\mathrm{d}V=\mu$. The zero function solves the smooth equation $\Omega^n=c\mu$ in the background $\Omega$, so its subsolution threshold is $b_\Omega(1)=c$, by \cref{prop:individual-normalizer-lp}. For $n\geq2$, the normalized $(n-1)$st power of the Gauduchon metric conformal to $\Omega$ is a smooth strictly positive maximizer of the dual problem. By \cref{cor:smooth-dual-exact-equation}, any two threshold subsolutions differ by a constant; comparing $w$ with zero shows that $w$ is constant. For $n=1$, the positive current $\Omega+\ddc w$ dominates its extended measure, hence dominates $c\mu=\Omega$, while both have total mass $\int_X\Omega=c$ by Stokes' theorem. Thus $\ddc w=0$, and $w$ is again constant.

\emph{Solutions in the ordinary type.} Finally, let $u$ be in the ordinary singularity type of $\sigma_D$ with $\mathrm{MA}^{\mathrm{ext}}_\theta(u)=c'\mu$ for some $c'\geq0$. Then $\nu(u,D_j)=a_j$, and the same reduction gives a bounded $w\in\PSH(X,\Omega)$ with $\mathrm{MA}_\Omega(w)=c'\mu$. Suppose $c'<c$. For every constant $C$, the bounded functions $w$ and $C$ satisfy $\mathrm{MA}_\Omega(w)=(c'/c)\mathrm{MA}_\Omega(C)$, so \cref{thm:bounded-strict-domination} gives $w\geq C$, which is impossible for a bounded $w$. The case $c'>c$ is symmetric, with the roles of $w$ and the constants exchanged. Therefore $c'=c$, and the rigidity just proved gives $u=u_0+\mathrm{constant}$.
\end{proof}


When $\theta$ is Hermitian and $n\geq2$, write $\mu=f\,\mathrm{d}V$ relative to a smooth strictly positive probability volume $\mathrm{d}V$. For every $u\in\PSH(X,\theta)$ with $\theta+\ddc u\geq[D]$ and $\mathrm{MA}^{\mathrm{ext}}_\theta(u)\geq cf\,\mathrm{d}V$, in particular for the potential $u_0$ of \cref{thm:divisorial-residual-equation} with $c=\mathcal V_\mu(\theta;D)$, \cref{prop:threshold-divisorial-constraint} bounds the coefficient $c$ in terms of the subsolution threshold $b_\theta(f)$ of \cref{thm:individual-density-subsolution} and the linear threshold of the class of the divisor. The residual problem prescribes the divisor, whereas the threshold $b_\theta(f)$ allows all $\theta$-plurisubharmonic functions.

\Needspace{12\baselineskip}
Along a family $D_t=tD$, $t\geq0$, in the feasible interval, the root of \eqref{eq:divisorial-gaud-volume} is concave and nonincreasing in $t$: every test quotient is affine in $t$ with nonpositive slope, and both properties pass to the infimum. On a subinterval where $\theta-t\eta_D$ has a smooth positive representative, choose a minimizing Gauduchon metric $\gamma_t$ for that background. \Cref{prop:gaud-volume-first-variation} then gives the sharper formula
\begin{equation}\label{eq:divisorial-volume-derivative}
 \frac{\mathrm{d}}{\mathrm{d}t}\mathcal V_\mu(\theta;tD)^{1/n}
       =-\frac{\deg_{\gamma_t}D}{D_\mu(\gamma_t)}<0.
\end{equation}
In dimension one the derivative is $-\deg D$, since the denominator is one. If the residual background at an endpoint of the interval is nef, \cref{thm:gaud-volume-nef-limit} also expresses its volume as the limit of the scalars of the smooth normalized equations for the perturbations by $\epsilon h$, with $h$ a fixed Hermitian form and $\epsilon\downarrow0$.

\Needspace{13\baselineskip}
\subsection{Varying divisorial coefficients}

For a fixed support, the residual equation depends smoothly on the coefficients as long as the residual background stays in $\mathscr P$. The smooth remainders then give convergence in capacity of the solutions, and they determine all generic Lelong numbers of the limit on every modification. Keep $D_j$, $s_j$ and $\eta_j$ as above, and put $\sigma_j=\log|s_j|^2$.

For $a=(a_1,\ldots,a_N)\in\mathbb R_{\geq0}^N$ put
\[
 D_a=\sum_j a_jD_j,\qquad
 \sigma_a=\sum_{a_j>0}a_j\sigma_j,\qquad
 \beta_a=\theta-\sum_j a_j\eta_j,
\]
where zero coefficients are omitted from the logarithmic sum. Let $\mathcal A$ be the set of such $a$ for which $\beta_a$ has a smooth positive representative.

\Needspace{14\baselineskip}
\begin{proposition}\label{prop:finite-divisorial-equation-limits}
There are unique smooth functions $v_a$, $a\in\mathcal A$, such that
\[
 \int_Xv_a\,\mathrm{d}\mu=0,\qquad
 \Omega_a=\beta_a+\ddc v_a>0,\qquad
 \Omega_a^n=c(a)\mu,
 \quad c(a)=\mathcal V_\mu(\beta_a)>0.
\]
The maps $a\mapsto v_a$ and $a\mapsto c(a)$ are smooth, and every smooth norm of $v_a$ is bounded on compact subsets of $\mathcal A$. Here smoothness at a parameter with a zero coefficient means that the family is the restriction of a smooth family defined on an open subset of $\mathbb R^N$.

Set $u_a=\sigma_a+v_a$. If $a^{(k)}\in\mathcal A$ and $a^{(k)}\to a\in\mathcal A$, then $u_{a^{(k)}}\to u_a$ in $L^1$ and in capacity, and smoothly on compact subsets of $X\setminus\bigcup_jD_j$. Moreover,
\[
 \bigl\|\mathrm{MA}^{\mathrm{ext}}_\theta(u_{a^{(k)}})
       -\mathrm{MA}^{\mathrm{ext}}_\theta(u_a)\bigr\|_{\mathrm{TV}}
       =|c(a^{(k)})-c(a)|\longrightarrow0,
\]
where $\|\cdot\|_{\mathrm{TV}}$ is the total variation norm. For every modification $\pi:Y\to X$ with $Y$ smooth and every prime divisor $E\subset Y$,
\[
 \nu(u_a\circ\pi,E)=\sum_j a_j\operatorname{ord}_E(\pi^*D_j),
\]
where $\operatorname{ord}_E(\pi^*D_j)$ is the vanishing order along $E$ of the pullback of a local equation of $D_j$. In particular, the limit $u_a$ has exactly the ordinary singularity type of $\sigma_a$, and no further singularity appears in this finite family.
\end{proposition}
\begin{proof}
The set of parameters in $\mathbb R^N$ for which $\beta_a$ has a smooth positive representative is open, and the smooth equation gives $v_a$ and $c(a)$ on it. The argument in the proof of \cref{prop:gaud-volume-first-variation} gives smooth dependence of both quantities on the parameter, with the fixed mean normalization; uniqueness makes the local families agree. Covering a compact set of parameters by finitely many such neighborhoods gives the asserted bounds.

The Poincar\'e--Lelong formula gives
\[
 \theta+\ddc u_a=[D_a]+\Omega_a,
\]
so $u_a\in\PSH(X,\theta)$. Its extended measure equals $\Omega_a^n=c(a)\mu$ off the union of the divisors, and both measures vanish on that union; this also covers the case $D_a=0$. Each fixed $\sigma_j$ is locally integrable, so the convergence of the coefficients and the smooth convergence of $v_a$ give convergence in $L^1$. Off the union of the divisors every $\sigma_j$ is smooth, which gives smooth convergence there. For convergence in capacity, remove a neighborhood of this analytic set of arbitrarily small Hermitian capacity; on its compact complement the convergence is uniform. Since the union of the divisors has capacity zero, this proves convergence in capacity. The identity for the total variation distances follows from the equations and $\mu(X)=1$.

Finally, after pullback the logarithmic potential $\sigma_j\circ\pi$ has generic Lelong number $\operatorname{ord}_E(\pi^*D_j)$ along $E$, whereas $v_a\circ\pi$ is smooth. Additivity of Lelong numbers for the finite sum proves the formula on every modification. Since $v_a$ is smooth also at the limiting parameter, $u_a$ has the ordinary singularity type of $\sigma_a$.
\end{proof}

The smooth bounds above use that the limiting parameter lies in $\mathcal A$. When the sections and weights move, \cref{thm:moving-divisorial-potential-stability} gives $L^1$ convergence of the potentials normalized by $\sup_X=0$; when the residual background is fixed, it also allows uniformly bounded measurable right-hand sides converging in $L^1$. On a modification, the pullback of a positive probability density can vanish along the exceptional divisor, so the smooth equation with positive density no longer applies directly. For a minimal exceptional pole, \cref{prop:exceptional-density-obstruction} gives a vanishing order of the density that is necessary for a smooth residual, but this order alone does not ensure smoothness.

Multiplier ideals need not stabilize in this family. At a smooth point of $D_j$ outside the other divisors, take a local coordinate $z_1$ defining $D_j$. There $u_a=a_j\log|z_1|^2+O(1)$, so
\[
 \mathcal I(u_a)=(z_1^{\lfloor a_j\rfloor})
\]
by the one-variable integral in \cref{ex:ideal-disc} and Fubini's theorem. Thus, if the parameter stays in $\mathcal A$ while $a_j$ increases to an integer from below, this local ideal jumps at the limit. Convergence of the Lelong numbers to their limiting values does not imply that the multiplier ideals are eventually equal to those of the limit.


\subsection{A ruled residual and its subtraction cost}
\label{subsec:ruled-residual-gauduchon-cost}

On the ruled surface of \cref{sec:nonclosed-ruled-trace-masses}, the Gauduchon variational volume has a different dependence on the data from the two volumes of an ordinary type. Keep $Y=\mathbb P^1\times\mathbb P^1$, $\theta=cw_1+aw_2$, $0\leq b<c$ and $\beta=(c-b)w_1+aw_2$. Let $m>0$ be smooth on the first curve, with $\int mw=1$, and put
\[
 \mu=mw_1\wedge w_2,\qquad k=c-b,\qquad H=\int_{\mathbb P^1}\frac{m}{a}w.
\]
In dimension two the Gauduchon tests are all smooth positive $(1,1)$-forms $T$ with $\ddc T=0$, and $D_\mu(T)=\int_Y\sqrt{T^2/\mu}\,\mu$.

\enlargethispage{4\baselineskip}
\begin{proposition}\label{prop:ruled-residual-gauduchon-volume}\index{Gauduchon variational volume!ruled residual}
The all-test infimum is attained and
\begin{equation}\label{eq:ruled-residual-gauduchon-volume}
 \mathcal V_\mu(\theta;bD)=\mathcal V_\mu(\beta)=\frac{2(c-b)}{H},
 \qquad \mathcal V_\mu(\theta)=\frac{2c}{H}.
\end{equation}
Writing $V=2k/H$ and $g=Vm/(2a)$, a minimizing test is the closed positive form
\[
 T=\frac ga w_1+w_2.
\]
\end{proposition}
\begin{proof}
Since $\int gw=k$, solve $\ddc f=(g-k)w$ on the first sphere. Then $\Omega=\beta+\ddc f=gw_1+aw_2>0$ and $\Omega^2=V\mu$. For every Gauduchon test, Stokes and the matrix arithmetic--geometric mean inequality give
\[
 \int_Y\beta\wedge T=\int_Y\Omega\wedge T
       \geq\sqrt V\,D_\mu(T).
\]
The displayed $T$ is closed, since $g/a$ depends only on the first factor, and $T=a^{-1}\Omega$, so equality holds pointwise. Explicitly,
\[
 \int_Y\beta\wedge T=2k,
 \qquad D_\mu(T)=\sqrt V\int\frac ma w=\sqrt V H.
\]
This proves the infimum over all tests. The curvature of $D$ is $w_1$, so the divisorial residual is indeed $\beta$. Replacing $k$ by $c$ gives the formula for $\theta$.
\end{proof}

Fix a real constant $d$ and a smooth real first-factor function $e$. On an interval where $c_s=c-sd>b$ and $a_s=a-se>0$, keep $b,m$ fixed and put
\[
 k_s=c_s-b,\qquad H_s=\int\frac{m}{a_s}w,
 \qquad V_s=\frac{2k_s}{H_s},\qquad
 T_s=\frac{V_sm}{2a_s^2}w_1+w_2.
\]
\begin{corollary}\label{cor:ruled-residual-subtraction-cost}
For the residual background $\beta_s=k_sw_1+a_sw_2$,
\begin{equation}\label{eq:ruled-residual-subtraction-cost}
 \frac{\mathrm d}{\mathrm ds}\sqrt{V_s}
     =-\frac{\int_Y(d w_1+e w_2)\wedge T_s}{D_\mu(T_s)}.
\end{equation}
Consequently, for every compact segment $[s_0,s_1]$ in that interval,
\[
 \sqrt{V_{s_0}}-\sqrt{V_{s_1}}
   =\int_{s_0}^{s_1}
       \frac{\int_Y(d w_1+e w_2)\wedge T_s}{D_\mu(T_s)}\,\mathrm ds.
\]
\end{corollary}
\begin{proof}
We have $H_s'=\int em/a_s^2\,w$, $D_\mu(T_s)=\sqrt{V_s}H_s$ and
\[
 \int_Y(d w_1+e w_2)\wedge T_s=d+\frac{V_s}{2}H_s'.
\]
Differentiating $V_s=2k_s/H_s$ gives \eqref{eq:ruled-residual-subtraction-cost}; integration proves the second identity. This is the explicit first variation of \cref{prop:gaud-volume-first-variation} along the residual family.
\end{proof}

For $c=1$, $b=1/4$, $a=2+\tau$ and $m=1$, the height $\tau$ has distribution $\mathrm d\tau/2$ under $w$. Hence $H=(\log3)/2$, and the residual and original variational volumes are $3/\log3$ and $4/\log3$. The ordinary-type volumes from \eqref{eq:nonclosed-ruled-ordinary-volumes} are instead $3/2$ and $9/2$, independent of $m$. The subtraction identity differentiates the density-dependent variational volume. Pure-pullback invariance of the ordinary-type masses does not identify an arbitrary smooth positive probability density upstairs with the degenerate pullback of $\mu$.


\section{A divisorial threshold on the Hopf threefold}
\label{sec:hopf-divisorial-volume}

On the Hopf threefold, the pullback of a hyperplane has an explicit divisorial threshold. The residual metric solves its normalized equation with zero potential, so the residual volume can be computed directly, also at the nef endpoint.

In this section $X$ is the Hopf threefold of \cref{chap:hopf}, with its projection $\pi\colon X\to\mathbb P^2$, and we write $\zeta$ for its vertical form, so that $\omega=\alpha+\zeta$. These forms satisfy
\[
 \alpha^3=0,\quad\zeta^2=0,\quad
 \ddc\zeta=-\alpha^2,\quad\int_X\alpha^2\wedge\zeta=2.
\]
Let $H$ be the inverse image under $\pi$ of a projective line in $\mathbb P^2$. The pulled-back Fubini--Study metric gives $\eta_H=\alpha$ and a logarithmic potential $\sigma_H$ with $\ddc\sigma_H=[H]-\alpha$.

\begin{proposition}\label{prop:hopf-divisorial-threshold}
The divisorial threshold is $\tau_\omega(H)=1$. For $\mu_0=\omega^3/6$ and $0\leq a\leq1$,
\[
 \mathcal V_{\mu_0}(\omega;aH)=6(1-a)^2.
\]
For $a<1$ this is the scalar of the equation solved, as in \cref{thm:divisorial-residual-equation}, by $a\sigma_H$ plus a smooth function. At $a=1$ the residual background is nef, and the logarithmic potential $\sigma_H$ has zero non-pluripolar Monge--Amp\`ere measure. For $a<1$, every residual background has infinite upper mass.
\end{proposition}
\begin{proof}
The residual background is
\[
 R_a=(1-a)\alpha+\zeta.
\]
It is positive for $0\leq a<1$ and semipositive for $a=1$, so every coefficient $a\leq1$ is feasible. For the reverse inequality, consider the metrics $\gamma_t=\alpha+t\zeta$, $t>0$. They are Gauduchon, since
\[
 \ddc(\gamma_t^2)=-2t\alpha^3=0.
\]
\Needspace{6\baselineskip}
Using $\zeta^2=0$, we obtain
\[
 \frac{\int_X\omega\wedge\gamma_t^2}{\deg_{\gamma_t}H}
       =\frac{1+2t}{2t}\longrightarrow1\qquad(t\to\infty),
\]
and the threshold formula gives $\tau_\omega(H)=1$.

For $a<1$,
\[
 R_a^3=3(1-a)^2\alpha^2\wedge\zeta
       =(1-a)^2\omega^3=6(1-a)^2\mu_0.
\]
Hence the zero potential solves the residual smooth equation, and uniqueness of the scalar gives the volume. At $a=1$, the form $R_1=\zeta$ is nef, because $\zeta+\epsilon\omega>0$ for $\epsilon>0$, and the same computation gives
\[
 \mathcal V_{\mu_0}(\zeta+\epsilon\omega)
       =6\epsilon^2(1+\epsilon)\longrightarrow0.
\]
By \cref{thm:gaud-volume-nef-limit}, the residual volume at $a=1$ is zero. Moreover $\omega+\ddc\sigma_H=[H]+\zeta$, so the non-pluripolar measure of $\sigma_H$ is $\zeta^3=0$.

Finally, fix $a<1$ and choose $0<\delta\leq\min\{1-a,1\}$, so that $R_a-\delta\omega\geq0$. For the smooth functions $V_j$ of \cref{thm:hopf-unbounded},
\[
 R_a+\ddc(\delta V_j)
       =(R_a-\delta\omega)+\delta(\omega+\ddc V_j)>0.
\]
Monotonicity of the determinant on positive matrices bounds the mass of this form below by $\delta^3\int_X(\omega+\ddc V_j)^3$, which tends to infinity. Thus a finite residual volume and an explicit divisorial threshold coexist with unbounded upper mass.
\end{proof}
